\documentclass{report}
\usepackage{amsmath,amssymb}
\setlength{\parindent}{0mm}
\setlength{\parskip}{1em}
\begin{document}
\begin{center}
\rule{6in}{1pt} \
{\large Dong-Lin Sun \\
{\bf A block GMRES method with deflated restarting for solving linear systems with multiple shifts and multiple right-hand sides}}

9747 AG Groningen \\ The Netherlands
\\
{\tt d.sun@rug.nl}\\
Ting-Zhu Huang\\
Yan-Fei Jing\\
Bruno Carpentieri\end{center}

We consider solutions of linear systems with multiple right-hand sides,
in which the coefficient matrices differ from each other by a scalar
multiple of the identity,
\begin{equation}\label{ABL}
(A-\sigma_{i}I)X_{i} = B,~\text{with} ~i = 1, 2, \ldots, L,
\end{equation}
where $A-\sigma_{i}I \in \mathbb C^{n\times n}$ are square nonsingular
matrices of large dimension $n$, $B = [b^{(1)}, b^{(2)},\ldots, b^{(p)}]$
is a full rank matrix consisting of $p\ll n$ given right-hand sides
$b^{(i)}$, and $\{\sigma _{i}\}_{i = 1}^{L}\subset \mathbb C$ are called
shifts. Many large-scale scientific applications, such as lattice quantum
chromo dynamics \cite{frommer1995} and PageRank problems
\cite{haveliwala2003topic}, to name but a few, require the solution of
sequences of linear systems with both multiple shifts and multiple
right-hand sides given simultaneously.

We propose a new shifted block Krylov subspace algorithm that solves all
of the given shifted linear systems simultaneously, and it has the
ability to dump the negative effect of small eigenvalues from the
convergence. Our block GMRES with deflated restarted for shifted linear
systems (BGMRES-DR-Sh) can be seen as an extension of the BGMRES-DR
method proposed by Morgan in~\cite{Morgan2005} to the case of shifted
linear systems.
This approach can be computationally more efficient than applying Krylov
subspace methods to the sequence of linear systems with multiple
shifts, or block Krylov subspace methods to each single shifted linear
system with multiple right-hand sides separately.

The key idea behind BGMRES-DR-Sh is to apply the block GMRES-DR algorithm
to the base shifted block system, and then to generate the approximate
solutions to the additional shifted block systems by imposing
collinearity with the computed base block residual. Moreover, our method
is enhanced with a seed selection strategy to handle the case of almost
linear dependence of the right-hand sides. After solving the seed block
system by the BGMRES-DR-Sh method, a Galerkin projection of the block
residual onto the shifted block Krylov subspace generated by the seed
block system is considered to obtain approximate solutions for the
additional block systems. The whole process is repeated until all of the
shifted block systems have converged.

In our talk we present the main lines of development of the BGMRES-DR-Sh
method, describing its theoretical properties, and we report on numerical
experiments on a set of representative matrix problems having size up to
1.5M unknowns arising from different fields. We analyse various aspects
such as performance comparison against other Krylov methods for solving
multi-shifted linear systems, the effect of the seed selection strategy,
and the sensitivity of our method to the accuracy of the
eigen-computation. Finally, we discuss how to combine it with a right
preconditioning technique that maintains the block shift-invariance
property.

{\small} \begin{thebibliography}{99}
\bibitem{frommer1995} A. Frommer, B. N{\"o}ckel, S. G{\"u}sken, T.
Lippert and K. Schilling, Many masses on one stroke: Economic computation
of quark propagators. \emph{Int. J. Mod. Phys. C}, \textbf{6} (1995)
627-638.
\bibitem{haveliwala2003topic} T. H. Haveliwala, Topic-sensitive pagerank:
A context-sensitive ranking algorithm for web search. \emph{IEEE Trans.
Knowl. Data Eng.}, \textbf{15} (2003) 784-796.
\bibitem{Morgan2005} R. B. Morgan, Restarted block-GMRES with deflation
of eigenvalues. \emph{Appl. Numer. Math}, \textbf{54} (2005) 222-236.
\end{thebibliography}


\end{document}
